\chapter{Gruppenwirkungen und innere Symmetrien}

\section{Wirkungen und Permutationsdarstellungen}

\begin{definition}[Gruppenwirkung und Äquivarianz]
\label{b40guide:def:action}
Eine \emph{linke Wirkung} einer Gruppe $G$ auf einer Menge $X$ ist eine
Abbildung $(g,x)\mapsto g\cdot x$ von $G\times X$ nach $X$ mit
\[
 1\cdot x=x,\qquad (gh)\cdot x=g\cdot(h\cdot x)
 \quad(g,h\in G,\ x\in X).
\]
Eine Menge mit einer solchen Wirkung heißt \emph{$G$-Menge}.
Für $G$-Mengen $X,Y$ heißt eine Abbildung $u:X\to Y$
\emph{$G$-äquivariant}, wenn $u(g\cdot x)=g\cdot u(x)$ gilt.
Eine äquivariante Bijektion heißt \emph{Isomorphismus von $G$-Mengen}.
\end{definition}

\begin{definition}[Permutationsdarstellung und Wirkungskern]
\label{b40guide:def:permutation-representation}
$\operatorname{Sym}(X)$ ist die Gruppe aller Permutationen von $X$
unter Komposition. Eine \emph{Permutationsdarstellung} von $G$ auf $X$
ist ein Homomorphismus $\rho:G\to\operatorname{Sym}(X)$.
Zu einer Wirkung gehört $\rho(g)(x)=g\cdot x$; ihr \emph{Kern} ist
\[
 K=\{g\in G:g\cdot x=x\text{ für alle }x\in X\}.
\]
Die Wirkung heißt \emph{treu}, wenn $K=\{1\}$, und \emph{frei}, wenn
aus $g\cdot x=x$ für $g\in G$ und $x\in X$ stets $g=1$ folgt.
Sie heißt \emph{transitiv}, wenn $X\ne\varnothing$ und zu je zwei
$x,y\in X$ ein $g\in G$ mit $g\cdot x=y$ existiert.
Eine freie transitive Wirkung heißt \emph{regulär}.
\end{definition}

\begin{theorem}[Wirkungen entsprechen Permutationsdarstellungen]
\label{b40guide:thm:action-representation}
Die Zuordnungen $\rho(g)(x)=g\cdot x$ und $g\cdot x=\rho(g)(x)$
stellen eine eindeutige Entsprechung zwischen Wirkungen von $G$ auf $X$
und Homomorphismen $G\to\operatorname{Sym}(X)$ her.
Dabei gilt $\rho(g)^{-1}=\rho(g^{-1})$.
Die Darstellung ist genau dann injektiv, wenn die Wirkung treu ist.
\cite{milne}
\end{theorem}

\begin{definition}[Bahn und Stabilisator]
\label{b40guide:def:orbit-stabilizer}
Für $x\in X$ heißen
\[
 G\cdot x=\{g\cdot x:g\in G\},\qquad
 G_x=\{g\in G:g\cdot x=x\}
\]
\emph{Bahn} beziehungsweise \emph{Stabilisator} von $x$.
\end{definition}

\begin{theorem}[Stabilisatoren und Wirkungskern]
\label{b40guide:thm:stabilizers-kernel}
Für jede Wirkung gilt $G_x\leq G$ und
\[
 G_{g\cdot x}=gG_xg^{-1},\qquad
 K=\bigcap_{x\in X}G_x\mathrel{\trianglelefteq}G.
\]
Für $X=\varnothing$ ist der Durchschnitt als $G$ zu lesen.
Die Formel $(gK)\cdot x=g\cdot x$ definiert eine treue Wirkung
von $G/K$ auf $X$ mit denselben Bahnen.
\cite{milne}
\end{theorem}

\begin{theorem}[Bahnen und Stabilisatorformel]
\label{b40guide:thm:orbit-stabilizer}
Die Bahnen sind die nichtleeren invarianten Teile der Zerlegung von $X$
in transitive $G$-Mengen. Für jedes $x\in X$ ist
\[
 G/G_x\longrightarrow G\cdot x,\qquad gG_x\longmapsto g\cdot x
\]
ein Isomorphismus von $G$-Mengen; links wirkt $G$ durch Linksmultiplikation.
Insbesondere gilt $|G\cdot x|=[G:G_x]$, auch für unendliche Bahnen.
Bei einer regulären Wirkung ist $g\mapsto g\cdot x$ für jedes $x$
ein Isomorphismus von der regulären $G$-Menge $G$ nach $X$.
\cite{milne}
\end{theorem}

\section{Nebenklassenwirkungen und Cayleys Satz}

\begin{definition}[Nebenklassenwirkung und Core]
\label{b40guide:def:core}
Für $H\leq G$ bezeichnet $G/H$ hier die Menge der Linksnebenklassen;
eine Quotientengruppenstruktur wird dabei nicht verlangt.
Die \emph{Nebenklassenwirkung} ist $g\cdot(aH)=(ga)H$.
Der \emph{Core} von $H$ in $G$ ist
\[
 \operatorname{Core}_G(H)=\bigcap_{g\in G}gHg^{-1}.
\]
\end{definition}

\begin{theorem}[Kern der Nebenklassenwirkung]
\label{b40guide:thm:coset-core}
Die Nebenklassenwirkung ist transitiv, der Stabilisator von $H$ ist $H$,
und ihr Kern ist $\operatorname{Core}_G(H)$.
Dieser Core ist der größte Normalteiler von $G$, der in $H$ enthalten ist.
Die Wirkung liefert eine Einbettung
\[
 G/\operatorname{Core}_G(H)\longrightarrow\operatorname{Sym}(G/H).
\]
Insbesondere ist die Nebenklassenwirkung genau dann treu, wenn
$\operatorname{Core}_G(H)=\{1\}$.
\cite{milne}
\end{theorem}

\begin{theorem}[Transitive Wirkungen und konjugierte Untergruppen]
\label{b40guide:thm:transitive-classification}
Jede transitive $G$-Menge ist zu $G/H$ für einen Stabilisator $H$
isomorph. Für $H,L\leq G$ gilt
\[
 G/H\cong G/L\text{ als $G$-Mengen}
 \quad\Longleftrightarrow\quad
 L=gHg^{-1}\text{ für ein }g\in G.
\]
\cite{milne}
\end{theorem}

\begin{theorem}[Cayleys Darstellungssatz]
\label{b40guide:thm:cayley}
Die Abbildung
\[
 \lambda:G\longrightarrow\operatorname{Sym}(G),\qquad
 \lambda(g)(x)=gx
\]
ist ein injektiver Gruppenhomomorphismus. Jede Gruppe ist somit
zu einer Untergruppe einer symmetrischen Gruppe isomorph.
Die zugrunde liegende Wirkung auf $G$ ist regulär.
\cite{milne}
\end{theorem}

\section{Konjugation, Zentralisatoren und Automorphismen}

\begin{definition}[Konjugation, Zentralisator, Zentrum und Normalisator]
\label{b40guide:def:conjugation}
Die \emph{Konjugation} ist $g\cdot x=gxg^{-1}$.
Für $A\subseteq G$, $x\in G$ und $H\leq G$ seien
\[
 \begin{aligned}
 C_G(A)&=\{g\in G:ga=ag\text{ für alle }a\in A\},\\
 C_G(x)&=C_G(\{x\}),\qquad Z(G)=C_G(G),\\
 N_G(H)&=\{g\in G:gHg^{-1}=H\}.
 \end{aligned}
\]
Diese heißen \emph{Zentralisator}, \emph{Zentrum} und \emph{Normalisator}.
Die Konjugationsklasse von $x$ ist
$x^G=\{gxg^{-1}:g\in G\}$.
\end{definition}

\begin{theorem}[Gruppenspezifische Folgen der Konjugationswirkung]
\label{b40guide:thm:conjugation-orbits}
Konjugation definiert eine Wirkung von $G$ auf sich selbst mit
Stabilisator $C_G(x)$ an $x$ und Kern $Z(G)$.
Insbesondere sind $C_G(A)$ Untergruppen und $Z(G)\trianglelefteq G$.
Es gilt
\[
 |x^G|=[G:C_G(x)],\qquad
 x^G=\{x\}\ \Longleftrightarrow\ x\in Z(G).
\]
Die Mengen $x^G$ zerlegen $G$ in Konjugationsklassen.
\cite{milne,conradconj}
\end{theorem}

\begin{definition}[Innere und äußere Automorphismen]
\label{b40guide:def:inner-outer}
$\operatorname{Aut}(G)$ bezeichnet die Gruppe aller Gruppenautomorphismen
von $G$ unter Komposition. Die Abbildungen
$c_g(x)=gxg^{-1}$ heißen \emph{innere Automorphismen}; ihre Gruppe ist
$\operatorname{Inn}(G)=\{c_g:g\in G\}$.
Die \emph{äußere Automorphismengruppe} ist
$\operatorname{Out}(G)=\operatorname{Aut}(G)/\operatorname{Inn}(G)$.
\end{definition}

\begin{theorem}[Innere Automorphismen und Zentrum]
\label{b40guide:thm:inner-center}
Die Abbildung $g\mapsto c_g$ ist ein Homomorphismus mit Kern $Z(G)$.
Ferner gilt
\[
 G/Z(G)\cong\operatorname{Inn}(G),\qquad
 \operatorname{Inn}(G)\trianglelefteq\operatorname{Aut}(G),
\]
und für $\alpha\in\operatorname{Aut}(G)$ ist
$\alpha\circ c_g\circ\alpha^{-1}=c_{\alpha(g)}$.
Damit ist die Gruppe $\operatorname{Out}(G)$ wohldefiniert.
\cite{milne}
\end{theorem}

\begin{theorem}[Normalisator und Automorphismen einer Untergruppe]
\label{b40guide:thm:normalizer-centralizer}
$N_G(H)$ ist die größte Untergruppe $L\leq G$, in der $H$ ein
Normalteiler ist. Konjugation definiert einen Homomorphismus
\[
 N_G(H)\longrightarrow\operatorname{Aut}(H),\qquad
 g\longmapsto (h\mapsto ghg^{-1}),
\]
mit Kern $C_G(H)$. Somit gilt
$C_G(H)\trianglelefteq N_G(H)$, und
$N_G(H)/C_G(H)$ ist zu einer Untergruppe von $\operatorname{Aut}(H)$
isomorph. Ferner ist $H\trianglelefteq G$ genau dann, wenn $N_G(H)=G$.
\cite{milne}
\end{theorem}
\begin{definition}[Charakteristische Untergruppe]
Eine Untergruppe \(H\leq G\) heißt \emph{charakteristisch in \(G\)},
wenn \(\alpha(H)=H\) für jeden Automorphismus \(\alpha\) von \(G\) gilt.
\end{definition}

\begin{theorem}[Charakteristische Untergruppen]
Jede charakteristische Untergruppe ist normal.
Sind \(K\) charakteristisch in \(H\) und \(H\) charakteristisch in \(G\),
so ist \(K\) charakteristisch in \(G\).
Ist \(K\) charakteristisch in \(H\) und \(H\trianglelefteq G\),
so ist \(K\trianglelefteq G\).
Das Zentrum \(Z(G)\) und die Kommutatoruntergruppe \(G'\) sind
charakteristisch in \(G\). \cite{milne}
\end{theorem}
